\documentclass[11pt]{article}
\usepackage[T1]{fontenc}
\usepackage[utf8]{inputenc}
\usepackage{lmodern}
\usepackage[margin=1in]{geometry}
\usepackage{amsmath,amssymb,amsthm}
\usepackage{microtype}
\usepackage{booktabs}
\usepackage[hidelinks]{hyperref}
\hypersetup{pdftitle={The shifted subtree-root annulus for trees of diameter at most three},
  pdfauthor={Xiaojian Zeng},pdfsubject={Subtree polynomial roots of stars and double stars}}
\newtheorem{theorem}{Theorem}
\newtheorem{lemma}[theorem]{Lemma}
\DeclareMathOperator{\diam}{diam}
\newcommand{\C}{\mathbb C}
\newcommand{\abs}[1]{\left|#1\right|}
\title{The shifted subtree-root annulus for trees\\of diameter at most three}
\author{Xiaojian Zeng\\[3pt]
\small PBC School of Finance, Tsinghua University, 100190\\
\small \href{mailto:xiaojianzeng@foxmail.com}{xiaojianzeng@foxmail.com}}
\date{}

\begin{document}
\maketitle
\begin{abstract}
Brown and Mol conjectured that every subtree root of an $n$-vertex tree,
$n\geq2$, lies in the annulus centred at $-1/2$ with inner radius $1/2$
and outer radius $1/2+(n-1)^{1/(n-1)}$.
We prove this assertion for all trees of diameter at most three.
For every double star with at least one leaf at each centre, all nonzero
roots lie strictly between the two circles. Together with the known star
formula, this gives the full boundary classification in this class:
zero is the sole inner-boundary root, and outer equality occurs precisely
for even-order stars, at the unique root $-1-(n-1)^{1/(n-1)}$.
The proof uses a convex endpoint estimate and Rouch\'e's theorem.
The conjecture for arbitrary trees remains unresolved.
\end{abstract}

\section{The annulus and the restricted result}

Let $T$ be a finite simple tree of order $n$. For $1\leq k\leq n$, let
$s_k(T)$ count the vertex subsets of size $k$ whose induced subgraphs are
connected, each subset counted once. Its \emph{subtree polynomial} is
\[
 \Phi_T(z)=\sum_{k=1}^n s_k(T)z^k.
\]
Enumeration by subtree order is part of the classical study of mean
subtree order initiated by Jamison~\cite{Jamison1983}.
We use the connected-induced-subgraph convention of
Brown and Mol~\cite{BrownMol2020}. Since $s_1(T)=n$, zero is always
a simple root of $\Phi_T$.

Brown and Mol proposed the following annulus bound
\cite[Conjecture~4.2 in the preprint version]{BrownMol2020}:
for every tree of order $n\geq2$ and every zero $z$ of $\Phi_T$,
\begin{equation}\label{eq:annulus}
 \frac12\leq\abs{z+\frac12}\leq\frac12+(n-1)^{1/(n-1)}.
\end{equation}
They reported verification for all trees through order~18.
Jin, Ma and Wang~\cite[Theorem~1.3]{JinMaWang2026} proved the
origin-centred bound $|z|\leq1+(n-1)^{1/(n-1)}$, including its equality
classification. Their Conjecture~5.1 retains~\eqref{eq:annulus}.
The origin-centred disk does not imply either of the two shifted bounds.

Write $D(a,b)$ for the double star obtained from an edge $uv$ by attaching
$a$ leaves to $u$ and $b$ leaves to $v$, where $a,b\geq1$ are integers.
These are exactly the trees of diameter three.

\begin{theorem}\label{thm:main}
Let $T$ be a finite simple tree with $n\geq2$ vertices and
$\diam(T)\leq3$. Every zero of $\Phi_T$ satisfies~\eqref{eq:annulus}.
The simple zero $z=0$ is the only root on the inner boundary.
A root lies on the outer boundary if and only if $n$ is even and
$T\cong K_{1,n-1}$; in that case the unique such root is
$z=-1-(n-1)^{1/(n-1)}$.
In particular, every nonzero subtree root of $D(a,b)$, for all integers
$a,b\geq1$, satisfies
\begin{equation}\label{eq:strict}
 \frac12<\abs{z+\frac12}<\frac12+(a+b+1)^{1/(a+b+1)}.
\end{equation}
\end{theorem}

The contribution is the proof for the complete two-parameter double-star
family, allowing arbitrary order and arbitrary imbalance between $a$ and $b$.
The star formula and the radii in~\eqref{eq:annulus} are prior results
and conjectured bounds, respectively. The substantive estimate below
controls $|z|$ and $|1+z|$ together on the translated circle.
For comparison, Wagner and Wang~\cite[Theorem~3]{WagnerWang2026}
obtain root bounds for graphs with minimum degree at least $\alpha n$,
for fixed $\alpha>0$ and sufficiently large $n$.
Their tree-subgraph convention agrees with ours when the ambient graph
is a tree, but the degree hypothesis does not cover an unbounded family
of double stars, whose minimum degree is one.

\section{A comparison on the translated circle}

\begin{lemma}\label{lem:endpoint}
Let $c>0$, $K=2c^2+2c+1$, and $k\geq1$ be real. For $c\leq t\leq c+1$,
\begin{equation}\label{eq:endpoint}
 \frac{t^{-k}}{\sqrt{K-t^2}}\leq\frac{c^{-k}}{c+1}.
\end{equation}
\end{lemma}
\begin{proof}
All denominators are positive, since $K-(c+1)^2=c^2$.
Set $f_k(t)=t^{-k}(K-t^2)^{-1/2}$. Direct differentiation gives
\[
 (\log f_k)''(t)=\frac{k}{t^2}+\frac{K+t^2}{(K-t^2)^2}>0.
\]
Consequently $f_k''=f_k\bigl(((\log f_k)')^2+(\log f_k)''\bigr)>0$.
Convexity bounds $f_k$ by the larger endpoint value, and
\[
 f_k(c)=\frac{c^{-k}}{c+1},\qquad
 \frac{f_k(c+1)}{f_k(c)}=\left(\frac{c}{c+1}\right)^{k-1}\leq1.
\]
This proves~\eqref{eq:endpoint}.
\end{proof}

The geometry behind the lemma is elementary. On
$|z+1/2|=c+1/2$, set $t=|z+1|$ and $s=|z|$. Then
\begin{equation}\label{eq:geometry}
 c\leq t\leq c+1,\qquad s^2+t^2=2c^2+2c+1=K,
 \qquad s=\sqrt{K-t^2}\geq c.
\end{equation}
The identity follows by expanding both squared moduli; the interval
follows by measuring the distance from $-1$ to that circle.
Thus Lemma~\ref{lem:endpoint} applies directly to $1/(s t^k)$.

\section{Proof of Theorem~\ref{thm:main}}

First consider $D(a,b)$ and put
\[
 q=a+b\geq2,\qquad w=1+z,\qquad c=(q+1)^{1/(q+1)}>1.
\]
A connected vertex subset either contains neither centre, contains $u$
but not $v$, contains $v$ but not $u$, or contains both centres.
In the first case it is a single leaf. The generating polynomials of
the four disjoint classes are $qz$, $zw^a$, $zw^b$, and $z^2w^q$.
Hence
\begin{equation}\label{eq:counting}
 \Phi_{D(a,b)}(z)=zQ(z),\qquad
 Q(z)=q+w^a+w^b+zw^q.
\end{equation}
The polynomial $Q$ is monic of degree $q+1$, and $Q(0)=q+2$.

\subsection*{The closed inner disk}
If $|z+1/2|\leq1/2$, then $|z|\leq1$, $|w|\leq1$ and
$\operatorname{Re}w\geq0$. For $q\geq4$,
\[
 |w^a+w^b+zw^q|\leq3<q,
\]
so $Q(z)\ne0$. For $q=3$ the same comparison is strict except possibly
at $z=0,-1$: away from these two points both $|z|$ and $|w|$ are
strictly less than one. At the exceptions, $Q(0)=q+2$ and $Q(-1)=q$.
For $q=2$, necessarily $a=b=1$, and
\[
 |2+2w|\geq\operatorname{Re}(2+2w)\geq2>1\geq|zw^2|.
\]
Thus $Q$ has no zero in the closed inner disk for every $a,b\geq1$.

\subsection*{The outer circle}
On $|z+1/2|=c+1/2$, use $s,t,K$ from~\eqref{eq:geometry}.
For fixed $t>1$, convexity of $x\mapsto t^x+t^{q-x}$ on $[1,q-1]$
gives
\[
 t^a+t^b\leq t+t^{q-1}.
\]
For $q=2$ this is an equality on a one-point interval.
The triangle inequality and Lemma~\ref{lem:endpoint} now imply
\begin{align}
 \frac{|q+w^a+w^b|}{|zw^q|}
 &\leq\frac{q+t^a+t^b}{s t^q}\notag\\
 &\leq f_{q-1}(t)+f_1(t)+qf_q(t)\notag\\
 &\leq\frac{c+c^{q-1}+q}{(c+1)c^q}=:F.\label{eq:domination}
\end{align}
Every index in this use of the lemma is at least one.
Using $c^{q+1}=q+1$, direct simplification yields
\begin{equation}\label{eq:margin}
 1-F=\frac{(c-1)(q+1-c^2)}{c(c+1)(q+1)}>0.
\end{equation}
Indeed $c>1$ and $c^2=(q+1)^{2/(q+1)}<q+1$ because $q+1\geq3$.
This is a positive margin for each parameter pair; no positive lower
bound independent of $q$ is needed.

It follows that $|q+w^a+w^b|<|zw^q|$ everywhere on the outer circle.
Rouch\'e's theorem, in its usual holomorphic form
\cite[Corollary~5.4.3]{Lebl2026}, gives the same number of zeros inside
the circle for $Q(z)$ and $z(z+1)^q$.
Both zeros $0,-1$ of the latter polynomial lie strictly inside, so the
number is $q+1$, counted with multiplicity. This is exactly the degree
of $Q$, and the strict comparison excludes boundary zeros.
All zeros of $Q$ therefore satisfy the strict outer bound. Together with
the inner-disk exclusion and~\eqref{eq:counting}, this proves
\eqref{eq:strict}.

\subsection*{Stars and the boundary classification}
A tree of diameter one or two is a star. To see that a tree of diameter
three is a double star, take a longest path $x,u,v,y$. Every component off
this path attaches to exactly one of its vertices. Attachment at $x$ or
$y$, or depth at least two from $u$ or $v$, would produce a path of
length at least four. All off-path vertices are therefore leaves at
$u$ or $v$, and $x,y$ ensure that both leaf populations are nonempty.

For the star $K_{1,m}$, where $m=n-1\geq1$, the known formula
\cite[Section~1]{BrownMol2020} is
\begin{equation}\label{eq:star}
 \Phi_{K_{1,m}}(z)=z\bigl((1+z)^m+m\bigr).
\end{equation}
Every nonzero root has $w=1+z$ satisfying $w^m=-m$ and
$|w|=d:=m^{1/m}$. If $m\geq2$, then $d>1$, and
\[
 \abs{z+\frac12}=\abs{w-\frac12}\geq d-\frac12>\frac12.
\]
For $m=1$ the sole nonzero root is $-2$, with shifted modulus $3/2$.
Also $|w-1/2|\leq d+1/2$, with equality exactly when $w=-d$.
This value solves $w^m=-m$ precisely when $m$ is odd, equivalently
when $n$ is even, and gives the unique boundary root $z=-1-d$.
Zero is simple by~\eqref{eq:star}. Combining the stars with the strict
double-star result proves every assertion of the theorem.\qed

\section{Exact checks and limits}

The preceding proof covers all parameters. Supplementary exact
computations checked the following finite ranges, with $a\leq b$ to
avoid repeating isomorphic double stars:
\begin{center}
\begin{tabular}{@{}lll@{}}
\toprule
Check & Parameter range & Number of cases\\
\midrule
Connected-subset coefficients & $2\leq a+b\leq10$ & $25$\\
Both root regions & $2\leq a+b\leq24$ & $144$\\
Endpoint identity~\eqref{eq:margin} & $2\leq q\leq64$ & $63$\\
\bottomrule
\end{tabular}
\end{center}
Each root check covers every root of the specified polynomial.
These ranges describe a double-star family check, not a census of trees.
The coefficient counts were also recomputed by a separate graph-traversal
implementation, and a separately written rational-arithmetic checker
verified the same $288$ root-region tests.

Here are the transformations and the certificate criterion, to make
the finite claims reproducible. For $d=\deg Q$, the inner check uses
\begin{equation}\label{eq:innertransform}
 I(u)=(2u)^d Q\left(\frac{1-u}{2u}\right).
\end{equation}
Its constant term is one. Its leading coefficient is checked to be
nonzero, excluding loss of a root at $z=-1/2$.
The transformed roots are $u=1/(2z+1)$, so $|u|<1$ is equivalent to
$|z+1/2|>1/2$. For the outer check choose
\[
 c_0=\frac{1}{10^4}\max\{k\in\mathbb Z_{\geq0}:
            k^{q+1}<(q+1)10^{4(q+1)}\},\qquad R=c_0+\frac12.
\]
Integer bisection determines $c_0$ exactly. Clearing rational
denominators in $Q(Ru-1/2)$ gives the outer polynomial, whose strict
unit-disk containment implies $|z+1/2|<R<c+1/2$.

The first checker uses the classical Schur reduction
\cite{Duffin1969}. For a real polynomial $P$ of degree $d$ with leading
coefficient $A$ and constant coefficient $B$, set
\[
 P^*(u)=u^dP(1/u),\qquad G(u)=\frac{AP(u)-BP^*(u)}{u}.
\]
If $|B|<|A|$ and all zeros of $G$ lie in $|u|<1$, then the same holds
for $P$. For completeness, a unit-circle zero of $P$ would also be a
zero of $P^*$, and hence of $G$, which is impossible. On that circle
$|P^*|=|P|$, so Rouch\'e's theorem compares $AP$ with $uG$.
The leading coefficient of $G$ is $A^2-B^2>0$ and its degree is $d-1$.
Recursive verification down to a nonzero constant therefore certifies
all roots using only integer arithmetic and positive-content divisions.

The second checker uses a half-plane coordinate. The map
$M(v)=(1+v)/(1-v)$ straightens the circular boundaries of a lens
in Zeng~\cite{Zeng2026Lens}.
Its signed conjugate gives the scalar coordinate
\[
 s=-M(-u)=\frac{u-1}{u+1},\qquad
 u=\frac{1+s}{1-s},\qquad
 \operatorname{Re}s=\frac{|u|^2-1}{|u+1|^2}.
\]
Thus $|u|<1$ corresponds to $\operatorname{Re}s<0$.
Disk and half-plane stability algorithms and their connection are
treated by Duffin~\cite{Duffin1969}.
For each transformed polynomial, the second checker formed
$(1-s)^dP((1+s)/(1-s))$, rejected any degree loss at $u=-1$, and
verified a strictly positive first column in its exact rational Routh
table. The code and complete initial integer polynomials are included
in the accompanying reproducibility dossier; no approximate root
solver is used for these certificates.

The theorem supplies no conclusion for diameter four or greater.
A bounded primary-source comparison completed on 13 September 2026
located no equivalent all-parameter double-star annulus theorem in the
inspected literature; it does not establish exhaustive priority.
The finite checks above are supplementary to the analytic proof and
do not resolve the arbitrary-tree conjecture.

\paragraph{Computational and AI assistance.}
OpenAI language models accessed through Codex assisted with proof
development, exact-check code, literature screening and manuscript
preparation. A separate model context examined the proof and produced
the second checking implementation. This separation does not imply
statistical independence or external peer review. The proof is given
in full above; the computations use Python integer and rational
arithmetic, with SymPy for supplementary symbolic identities.

\bibliographystyle{plain}
\bibliography{references}
\end{document}
